\documentclass[../../main.tex]{subfiles}

\begin{document}

\chapter{Prosedur dan Manajemen Stack}
\label{chap:prosedur}

% =========================================================
% OUTCOME PEMBELAJARAN KHUSUS BAB
% =========================================================

\begin{subcpmkbox}
\SubCPMK{3.4} --- Mengembangkan prosedur/subrutin dengan mekanisme
pengiriman parameter (\eng{parameter passing}) dan mengelola
\eng{stack frame} melalui instruksi \asm{PUSH} dan \asm{POP},
register \reg{BP}, serta konvensi pemanggilan yang benar.

\vspace{0.4em}
\textbf{CPMK yang didukung:} \CPMK{03} \quad\textbullet\quad
\textbf{CPL yang didukung:} \CPL{03}
\end{subcpmkbox}

\begin{tujuanbox}
Setelah mempelajari bab ini, mahasiswa mampu:
\begin{enumerate}
    \item menjelaskan alasan pemakaian prosedur: modularitas,
          penggunaan ulang, dan abstraksi, beserta biaya
          pemanggilannya;
    \item menjelaskan kerja instruksi \asm{CALL} dan \asm{RET} serta
          membedakan pemanggilan \eng{near} dan \eng{far};
    \item menelusuri isi stack sel per sel pada setiap operasi
          \asm{PUSH}, \asm{POP}, dan \asm{CALL};
    \item membangun \eng{stack frame} standar melalui \eng{prologue}
          dan \eng{epilogue} yang benar;
    \item mengirim parameter melalui register, melalui stack, dan
          melalui variabel global beserta kelebihan serta kekurangan
          masing-masing cara;
    \item mengembalikan nilai prosedur melalui register \reg{AX} dan
          register \reg{DX}:\reg{AX} serta menjelaskan peran
          \asm{RET} \code{n};
    \item menulis prosedur rekursif untuk faktorial dan Fibonacci
          serta menjelaskan risiko \eng{stack overflow};
    \item mengenali dan memperbaiki kesalahan pengelolaan stack yang
          umum terjadi.
\end{enumerate}
\end{tujuanbox}

\begin{materibox}
\begin{enumerate}
    \item Alasan prosedur: modularitas, penggunaan ulang, abstraksi,
          dan biaya pemanggilan prosedur.
    \item Instruksi \asm{CALL} dan \asm{RET}; alamat kembali
          \reg{CS}:\reg{IP}; pemanggilan \eng{near} dan \eng{far}.
    \item Operasi stack: \asm{PUSH}, \asm{POP}, \asm{PUSHF},
          \asm{POPF}; register \reg{SP} dan \reg{BP}; arah tumbuh
          stack ke alamat rendah.
    \item \eng{Stack frame} standar: \eng{prologue}, akses parameter
          dengan \code{[BP+4]} dan \code{[BP+6]}, serta
          \eng{epilogue}.
    \item Konvensi pemanggilan dan tanggung jawab pembersihan stack.
    \item Nilai balik prosedur dan peran \asm{RET} \code{n}.
    \item Rekursi pada assembly: faktorial dan Fibonacci; kedalaman
          stack dan risiko \eng{stack overflow}.
    \item Kesalahan umum pengelolaan stack.
\end{enumerate}
\end{materibox}

\subfile{section-01}
\subfile{section-02}
\subfile{section-03}

\section{Contoh}
\label{sec:contoh-bab13}

Dua contoh berikut memakai percabangan dan perulangan pada program
utuh. Contoh pertama mencari nilai maksimum dengan \asm{LOOP} dan
\asm{JAE}. Contoh kedua mengurutkan array dengan perulangan
bersarang dan \asm{JBE}.

\subsection{Contoh 1: Mencari Nilai Maksimum}

Program pada Kode~\ref{lst:loop-maksimum} memeriksa setiap elemen
array dan menyimpan nilai terbesar pada \reg{AX}.

\lstinputlisting[
    language=MASM,
    caption={Program \file{loop-maksimum.asm} mencari nilai
    maksimum array.},
    label={lst:loop-maksimum}
]{\CodeDir/loop-maksimum.asm}

Alur percabangannya:

\begin{itemize}
    \item \lstasm|LEA SI, deret| menunjuk elemen pertama.
    \item \lstasm|MOV CX, BANYAK - 1| menyiapkan cacah perbandingan.
    \item \lstasm|CMP AX, [SI]| membandingkan kandidat dengan elemen
          saat ini.
    \item \lstasm|JAE LANJUT| melompat bila kandidat lebih besar atau
          sama. Lompatan ini memakai tafsiran tidak bertanda.
    \item \lstasm|MOV AX, [SI]| mengganti kandidat bila elemen lebih
          besar.
    \item \lstasm|LOOP CARI| mengulang sampai semua elemen diperiksa.
\end{itemize}

Keluaran program tampak pada Kode~\ref{lst:keluaran-maksimum}.

\begin{lstlisting}[
    style=sessiongaya,
    caption={Keluaran program \file{loop-maksimum.asm}.},
    label={lst:keluaran-maksimum}
]
Nilai maksimum: 93
\end{lstlisting}

\subsection{Contoh 2: Pengurutan Gelembung}

Program pada Kode~\ref{lst:bubble-sort} mengurutkan array byte
secara menaik dengan algoritme pengurutan gelembung.

\lstinputlisting[
    language=MASM,
    caption={Program \file{bubble-sort.asm} mengurutkan array
    dengan perulangan bersarang.},
    label={lst:bubble-sort}
]{\CodeDir/bubble-sort.asm}

Bagian penting program:

\begin{itemize}
    \item Perulangan luar memakai \reg{CX} untuk menentukan banyak
          lintasan.
    \item Pencacah luar disimpan dengan \asm{PUSH CX} sebelum
          perulangan dalam memakai \reg{CX}.
    \item Perulangan dalam membandingkan pasangan elemen yang
          berdekatan dengan \asm{CMP} dan \asm{JBE}.
    \item Bila elemen kiri lebih besar dari elemen kanan, kedua elemen
          ditukar.
    \item \asm{POP CX} memulihkan pencacah luar sebelum
          \asm{LOOP LINTASAN}.
\end{itemize}

Keluaran program tampak pada Kode~\ref{lst:keluaran-bubble}.

\begin{lstlisting}[
    style=sessiongaya,
    caption={Keluaran program \file{bubble-sort.asm}.},
    label={lst:keluaran-bubble}
]
Hasil pengurutan: 1 2 3 4 6 7 8 9
\end{lstlisting}

\begin{infobox}
\vspace{1.1em}
Kedua contoh memakai prosedur pembantu dan \eng{stack} untuk
pencetakan serta penyimpanan pencacah. Mekanisme \asm{CALL},
\asm{RET}, \asm{PUSH}, dan \asm{POP} dibahas pada bab berikutnya.
Pada bab ini, perhatian diarahkan pada alur \asm{CMP}, \asm{Jcc},
dan \asm{LOOP}.
\end{infobox}

\section{Praktikum}
\label{sec:praktikum-bab13}

\begin{praktikumbox}
\vspace{1.1em}
\textbf{Judul:} Menelusuri percabangan dan perulangan pada program
array.\\[0.3em]
\textbf{Sub-CPMK:} 3.4 \quad\textbullet\quad
\textbf{Alokasi waktu:} 100 menit \quad\textbullet\quad
\textbf{Bentuk:} individu
\end{praktikumbox}

\subsection{Tujuan Praktikum}

Setelah menyelesaikan praktikum ini, mahasiswa mampu menelusuri
percabangan dengan \asm{CMP} dan \asm{Jcc}, perulangan dengan
\asm{LOOP}, serta perulangan bersarang pada pengolahan array.

\subsection{Alat dan Bahan}

\begin{itemize}
    \item Komputer dengan emu8086 terpasang.
    \item DOSBox dengan TASM, TLINK, dan DEBUG.
    \item Berkas \file{loop-maksimum.asm} dan
          \file{bubble-sort.asm} dari \file{code/asm/}.
    \item Lembar pengamatan register dan \eng{flag}.
\end{itemize}

\subsection{Langkah Kerja}

\begin{enumerate}
    \item \textbf{Menguji pencarian maksimum.}
    \begin{enumerate}
        \item Rakit \file{loop-maksimum.asm} dengan TASM, lalu link
              dengan TLINK.
        \item Jalankan program dan catat keluarannya.
        \item Buka program pada emu8086 dan jalankan
              \eng{single step} sampai label \code{CARI}.
        \item Catat \reg{AX}, \reg{SI}, \reg{CX}, dan \flag{CF}
              setiap kali \lstasm|CMP AX, [SI]| selesai.
        \item Jelaskan kapan \lstasm|JAE LANJUT| diambil dan kapan
              tidak.
    \end{enumerate}

    \item \textbf{Mengubah data maksimum.}
    \begin{enumerate}
        \item Ubah array menjadi \code{25, 7, 120, 41, 93, 18, 60, 5}.
        \item Rakit ulang dan jalankan program.
        \item Pastikan nilai maksimum berubah menjadi 120.
    \end{enumerate}

    \item \textbf{Menguji pengurutan gelembung.}
    \begin{enumerate}
        \item Rakit dan jalankan \file{bubble-sort.asm}.
        \item Catat keluarannya.
        \item Telusuri satu lintasan luar dan catat nilai \reg{CX},
              \reg{SI}, dan \reg{DI}.
        \item Jelaskan peran \lstasm|JBE TIDAK_TUKAR| pada
              pembandingan.
    \end{enumerate}

    \item \textbf{Mengubah data pengurutan.}
    \begin{enumerate}
        \item Ubah array menjadi \code{9, 8, 7, 6, 5, 4, 3, 2}.
        \item Rakit ulang dan jalankan program.
        \item Pastikan hasilnya terurut menaik.
    \end{enumerate}

    \item \textbf{Membandingkan LOOP dan DEC/JNZ.}
    \begin{enumerate}
        \item Ganti \lstasm|LOOP BANDING| dengan pasangan
              \lstasm|DEC CX| dan \lstasm|JNZ BANDING|.
        \item Rakit ulang dan pastikan hasil pengurutan tetap sama.
        \item Jelaskan perbedaan penggunaan \reg{CX} pada kedua
              bentuk.
    \end{enumerate}

    \item \textbf{Menulis percabangan if--else.}
    \begin{enumerate}
        \item Tulis potongan program yang membandingkan \reg{AX} dan
              \reg{BX}.
        \item Bila \reg{AX} lebih besar, isi \code{hasil} dengan 1;
              bila tidak, isi dengan 0.
        \item Gunakan \asm{CMP}, \asm{JLE}, dan \asm{JMP}.
    \end{enumerate}
\end{enumerate}

\subsection{Hasil yang Dikumpulkan}

\begin{enumerate}
    \item Berkas \file{loop-maksimum.asm} dan
          \file{bubble-sort.asm} versi modifikasi.
    \item Tangkapan layar keluaran kedua program.
    \item Tabel pengamatan \reg{AX}, \reg{SI}, \reg{CX}, dan
          \eng{flag}.
    \item Penjelasan pemilihan \asm{JAE} dan \asm{JBE}.
    \item Potongan program \code{if--else}.
\end{enumerate}

\begin{asesmenbox}
\vspace{1.1em}
\textbf{Asesmen praktikum 11 (bobot 5\% dari komponen praktikum).}

\begin{tabular}{@{}p{6.2cm}cp{4.6cm}@{}}
    \toprule
    \textbf{Indikator} & \textbf{Bobot} & \textbf{Kriteria} \\
    \midrule
    Ketepatan penelusuran percabangan & 30\% &
    Keputusan \asm{JAE}/\asm{JBE} dijelaskan dengan benar \\
    \addlinespace
    Ketepatan penelusuran perulangan & 30\% &
    Nilai \reg{CX}, \reg{SI}, dan \reg{DI} tercatat benar \\
    \addlinespace
    Program modifikasi berjalan & 25\% &
    Keluaran berubah sesuai data baru \\
    \addlinespace
    Ketepatan potongan if--else & 15\% &
    Alur dan label disusun dengan benar \\
    \bottomrule
\end{tabular}
\end{asesmenbox}

\section{Latihan}
\label{sec:latihan-bab13}

\begin{enumerate}
    \item \textbf{Konseptual.} Sebutkan tiga alasan mengapa program
          sebaiknya dipecah menjadi prosedur. Jelaskan pula mengapa
          pemakaian prosedur tidak selalu membuat program menjadi
          lebih cepat, dan kapan prosedur sebaiknya dihindari.

    \item \textbf{Konseptual.} Jelaskan apa yang dikerjakan
          instruksi \asm{CALL} dan \asm{RET}. Mengapa \asm{JMP}
          tidak dapat dipakai sebagai pengganti \asm{CALL}?

    \item \textbf{Analisis stack.} Diketahui \reg{SS} $=$ \hex{2000}
          dan \reg{SP} $=$ \hex{FFFC}. Tentukan alamat fisik sel
          teratas stack. Kemudian tentukan nilai \reg{SP} dan isi
          sel tersebut setelah urutan instruksi berikut dijalankan:

\begin{lstlisting}[language=MASM, numbers=none]
    MOV AX, 1234h
    PUSH AX
    PUSH BX
    POP  CX
\end{lstlisting}

          \begin{enumerate}
              \item Berapa byte data yang masih tertinggal di stack
                    setelah instruksi terakhir?
              \item Apakah nilai \reg{BX} masih tersimpan di memori
                    stack? Jelaskan.
          \end{enumerate}

    \item \textbf{Analisis \eng{stack frame}.} Perhatikan potongan
          prosedur berikut.

\begin{lstlisting}[language=MASM, numbers=none]
Jumlah PROC
    PUSH BP
    MOV  BP, SP
    MOV  AX, [BP+4]
    ADD  AX, [BP+6]
    POP  BP
    RET  4
Jumlah ENDP
\end{lstlisting}

          \begin{enumerate}
              \item Tuliskan urutan instruksi pemanggil yang tepat
                    agar prosedur ini menghasilkan nilai benar.
              \item Nilai apa yang berada pada \code{[BP+8]} saat
                    prosedur berjalan?
              \item Mengapa pemanggil tidak boleh menambahkan
                    \lstasm|ADD SP, 4| setelah \asm{CALL}?
          \end{enumerate}

    \item \textbf{Analisis kesalahan.} Sebuah prosedur memakai
          register \reg{SI} dan \reg{DI} untuk mengolah data, tetapi
          hanya menyimpan \reg{BP} pada \eng{prologue}-nya. Setelah
          prosedur dipanggil, program pemanggil menghasilkan keluaran
          yang salah. Jelaskan penyebabnya dan tuliskan perbaikan
          yang diperlukan, dengan mengacu pada
          Tabel~\ref{tab:register-simpan}.

    \item \textbf{Penulisan program.} Tulislah prosedur
          \code{HitungRataRata} yang menerima dua parameter 16-bit
          melalui stack dan mengembalikan hasil bagi keduanya
          (dibulatkan ke bawah) pada register \reg{AX}. Sertakan
          \eng{prologue}, \eng{epilogue}, dan potongan program
          pemanggilnya.

    \item \textbf{Penulisan program.} Tulislah prosedur rekursif
          \code{Pangkat} yang menghitung $a^b$ untuk $b \ge 0$,
          dengan parameter $a$ dan $b$ dikirim melalui stack serta
          hasil dikembalikan pada \reg{AX}. Tentukan pula kasus
          basisnya dan jelaskan mengapa kasus basis itu diperlukan.

    \item \textbf{Perbandingan.} Bandingkan pemakaian stack pada
          prosedur faktorial rekursif dan faktorial iteratif untuk
          $n = \dec{10}$. Hitung perkiraan jumlah byte stack yang
          dipakai masing-masing versi, lalu simpulkan versi mana yang
          lebih tepat dipakai dan pada keadaan apa rekursi tetap
          bermanfaat.
\end{enumerate}

\begin{refleksibox}
Jawab pertanyaan berikut pada jurnal praktikum Anda.
\begin{enumerate}
    \item Bagian mana dari bab ini yang paling sulit dipahami:
          arah tumbuh stack, susunan \eng{stack frame}, atau konvensi
          pemanggilan? Apa yang Anda lakukan untuk mengatasinya?
    \item Setelah memeriksa isi stack pada \eng{debugger}, apakah
          gambaran Anda tentang ``pemanggilan prosedur'' berubah?
          Jelaskan bagian mana yang paling mengejutkan.
    \item Konsep apa yang masih perlu Anda latih sebelum masuk ke
          Bab~\ref{chap:interupsi}, yang juga memakai stack untuk
          menyimpan \reg{CS}, \reg{IP}, dan \reg{FLAGS}?
\end{enumerate}
\end{refleksibox}

\section{Rangkuman}
\label{sec:rangkuman-bab13}

Pada bab ini telah dibahas:

\begin{itemize}
    \item Label adalah nama simbolik untuk lokasi di dalam segmen
          kode. Label didefinisikan dengan titik dua dan ditulis
          tanpa titik dua saat menjadi tujuan lompatan.

    \item \asm{JMP} memindahkan kendali tanpa syarat, sedangkan
          \asm{Jcc} memindahkan kendali hanya bila syarat
          \eng{flag} terpenuhi.

    \item \asm{CMP} mengurangkan operand kedua dari operand pertama
          tanpa menyimpan hasil. Yang berubah hanya \eng{flag}.

    \item Lompatan bersyarat untuk bilangan bertanda memakai
          \asm{JG}, \asm{JGE}, \asm{JL}, dan \asm{JLE}. Untuk
          bilangan tidak bertanda dipakai \asm{JA}, \asm{JAE},
          \asm{JB}, dan \asm{JBE}.

    \item Pola \code{if} diterjemahkan menjadi \asm{CMP} diikuti
          satu lompatan bersyarat. Pola \code{if--else} memerlukan
          tambahan \asm{JMP} untuk melewati blok \code{else}.

    \item \asm{LOOP} mengurangi \reg{CX} dengan satu, lalu melompat
          bila \reg{CX} belum nol. Badan perulangan selalu dijalankan
          paling sedikit satu kali.

    \item Perulangan \code{while} memeriksa syarat sebelum badan,
          sedangkan \code{for} dengan cacah tetap dapat memakai
          \asm{LOOP}.

    \item Perulangan bersarang memerlukan penyimpanan pencacah luar,
          biasanya melalui \eng{stack} dengan \asm{PUSH} dan
          \asm{POP}.
\end{itemize}

\begin{checklistbox}
\vspace{1.1em}
Tandai setiap butir yang sudah Anda kuasai.

\begin{itemize}
    \item[\checkbox] Saya dapat menulis label dan menggunakannya
          sebagai tujuan lompatan.
    \item[\checkbox] Saya dapat menggunakan \asm{JMP}.
    \item[\checkbox] Saya dapat menggunakan \asm{CMP} untuk
          menyiapkan \eng{flag}.
    \item[\checkbox] Saya dapat memilih \asm{Jcc} untuk bilangan
          bertanda dan tidak bertanda.
    \item[\checkbox] Saya dapat menyusun pola \code{if} dan
          \code{if--else}.
    \item[\checkbox] Saya dapat menyusun perulangan dengan
          \asm{LOOP}.
    \item[\checkbox] Saya dapat menyusun perulangan bersarang untuk
          array.
\end{itemize}

\vspace{0.4em}
Bila masih ada butir yang belum bertanda, ulangi bagian yang
bersesuaian sebelum melanjutkan ke bab berikutnya.
\end{checklistbox}

\begin{infobox}
\vspace{1.1em}
\textbf{Rujukan bab.} Materi pada bab ini diringkas dari
\cite{gunawan2019}, Bagian 11, hlm.~101--114, dan
\cite{sto2001}, BAB XIII, hlm.~68--72.
\end{infobox}

\section{Evaluasi}
\label{sec:evaluasi-bab13}

\subsection{Soal Pilihan Ganda}

\begin{enumerate}
    \item Instruksi \asm{CALL} \eng{near} menyimpan \dots ke dalam
          stack.
          \begin{enumerate}
              \item register \reg{CS} dan \reg{IP}
              \item register \reg{IP} saja
              \item seluruh register umum
              \item register \reg{FLAGS} dan \reg{BP}
          \end{enumerate}

    \item Pada 8086, instruksi \asm{PUSH} bekerja dengan cara \dots
          \begin{enumerate}
              \item menambah \reg{SP} sebesar \dec{2}, lalu menyimpan
                    data
              \item mengurangi \reg{SP} sebesar \dec{2}, lalu
                    menyimpan data
              \item menyimpan data, lalu menambah \reg{SP} sebesar
                    \dec{2}
              \item menyimpan data, lalu mengurangi \reg{SP} sebesar
                    \dec{2}
          \end{enumerate}

    \item Setelah \eng{prologue} prosedur dijalankan, parameter
          pertama yang didorong paling akhir oleh pemanggil berada
          pada alamat \dots
          \begin{enumerate}
              \item \code{[BP+0]} \item \code{[BP+2]}
              \item \code{[BP+4]} \item \code{[BP+6]}
          \end{enumerate}

    \item Prosedur yang memakai \asm{RET} \code{4} berarti \dots
          \begin{enumerate}
              \item pemanggil harus menambahkan \lstasm|ADD SP, 4|
              \item prosedur membuang \dec{4} byte argumen setelah
                    mengambil alamat kembali
              \item prosedur mengembalikan nilai \dec{4} bit
              \item prosedur menyimpan empat register ke stack
          \end{enumerate}

    \item Pernyataan yang benar mengenai register \eng{callee-saved}
          pada konvensi 8086 adalah \dots
          \begin{enumerate}
              \item \reg{AX} dan \reg{DX} wajib dipulihkan prosedur
              \item \reg{BP}, \reg{SI}, dan \reg{DI} wajib dipulihkan
                    prosedur
              \item seluruh register segmen boleh dirusak prosedur
              \item \reg{SP} boleh berubah setelah \asm{RET}
          \end{enumerate}
\end{enumerate}

\subsection{Soal Uraian}

\begin{enumerate}
    \item Diketahui \reg{SS} $=$ \hex{3000} dan \reg{SP} $=$
          \hex{FFF8}. Gambarkan isi stack sel per sel (minimal lima
          sel) untuk urutan berikut: \lstasm|PUSH 000Ah|,
          \lstasm|PUSH 0014h|, \lstasm|CALL Hitung|,
          \lstasm|PUSH BP|, \lstasm|MOV BP, SP|. Tuliskan alamat
          setiap sel, isinya, dan nilai \reg{SP} serta \reg{BP}
          setelah setiap langkah.

    \item Tuliskan prosedur \code{HitungLuas} lengkap dengan
          \eng{prologue} dan \eng{epilogue}, lalu tuliskan pula
          bagian pemanggilnya. Jelaskan mengapa urutan \asm{PUSH}
          pada pemanggil menentukan nomor offset \code{[BP+4]} dan
          \code{[BP+6]}.

    \item Sebuah program berprosedur menghasilkan perilaku yang tidak
          menentu: kadang berjalan benar, kadang melompat ke alamat
          acak. Sebutkan empat kemungkinan penyebab yang berkaitan
          dengan pengelolaan stack, lalu jelaskan langkah pemeriksaan
          untuk masing-masing penyebab dengan bantuan
          \eng{debugger}.

    \item Bandingkan prosedur faktorial rekursif dan iteratif untuk
          $n = \dec{12}$ dari segi pemakaian stack, jumlah
          pemanggilan, dan risiko \eng{stack overflow}. Simpulkan
          kapan rekursi tetap merupakan pilihan yang tepat.
\end{enumerate}

\subsection{Kisi-kisi Penilaian}

\begin{table}[H]
    \centering
    \caption{Kisi-kisi evaluasi Bab~\ref{chap:prosedur}.}
    \label{tab:kisi-bab13}
    \small
    \begin{tabular}{@{}p{1.5cm}p{5.2cm}p{2.8cm}c@{}}
        \toprule
        \textbf{Sub-CPMK} & \textbf{Indikator yang diukur} &
        \textbf{Bentuk soal} & \textbf{Bobot} \\
        \midrule
        3.4 & Menjelaskan kerja \asm{CALL} dan \asm{RET} serta isi
        stack saat pemanggilan & Pilihan ganda 1, 2 & 15\% \\
        \addlinespace
        3.4 & Menganalisis \eng{stack frame} dan akses parameter
        melalui \reg{BP} & Pilihan ganda 3, 4; Uraian 2 & 30\% \\
        \addlinespace
        3.4 & Mengelola konvensi pemanggilan dan penyimpanan register
        & Pilihan ganda 5; Uraian 3 & 25\% \\
        \addlinespace
        3.4 & Menelusuri isi stack dan menilai pemakaian rekursi &
        Uraian 1, 4 & 30\% \\
        \bottomrule
    \end{tabular}
\end{table}

\begin{checklistbox}
\textbf{Checklist Pencapaian Kompetensi Bab~\ref{chap:prosedur}}

Seorang mahasiswa dinyatakan kompeten pada Sub-CPMK 3.4 apabila:
\begin{itemize}
    \item[\checked] dapat menjelaskan mekanisme \asm{CALL} dan
          \asm{RET} beserta isi stack yang terlibat;
    \item[\checked] dapat menggambar isi stack sel per sel pada
          setiap langkah pemanggilan prosedur;
    \item[\checked] dapat menulis prosedur berparameter dengan
          \eng{prologue}, \eng{epilogue}, dan \asm{RET} \code{n}
          yang benar;
    \item[\checked] dapat mengembalikan nilai prosedur melalui
          \reg{AX} serta menjelaskan pemakaian \reg{DX}:\reg{AX};
    \item[\checked] dapat menulis prosedur rekursif dengan kasus
          basis dan menjelaskan risiko \eng{stack overflow};
    \item[\checked] dapat menemukan dan memperbaiki kesalahan stack
          yang tidak seimbang pada program nyata.
\end{itemize}
\end{checklistbox}


\end{document}
